\phantomsection\addcontentsline{toc}{chapter}{Suplemento: propiedades y demostraciones}\label{mod:sup}
\chapter{Distribuciones finitas y secuencias}
\section{El símplex y el soporte}
Sea $V=\{1,\ldots,K\}$. El símplex es
$\Delta_{K-1}=\{p\in\mathbb R^K:p_i\ge0,\sum_i p_i=1\}$.
Su interior relativo exige $p_i>0$. El soporte es
$\operatorname{supp}(p)=\{i:p_i>0\}$. La convención
$0\log0=0$ está motivada por $\lim_{x\downarrow0}x\log x=0$.
Cuando $q_i=0$ y $p_i>0$, $-p_i\log q_i=+\infty$.
\section{Regla de la cadena con cuidado del soporte}
\textbf{Proposición.} Si $p(x_{1:T})>0$, entonces
\[p(x_{1:T})=p(x_1)\prod_{t=2}^T p(x_t\mid x_{1:t-1}).\]
\textit{Demostración.} Para todo prefijo del soporte, la definición de
condicional permite escribir
$p(x_{1:t})=p(x_{1:t-1})p(x_t\mid x_{1:t-1})$.
Al iterar, los factores de prefijos se cancelan telescópicamente.\hfill$\square$
Para prefijos de probabilidad cero, las condicionales pueden elegirse por
convención; la representación de la ley conjunta sigue siendo válida donde
está definida. Esto evita atribuir contenido probabilístico a expresiones
condicionales que no son identificables por datos.
\section{Markov no es una consecuencia de la factorización}
La hipótesis de orden $m$ afirma
$X_t\perp X_{1:t-m-1}\mid X_{t-m:t-1}$ bajo las condiciones de la ley.
La notación de independencia condicional es una propiedad de la distribución;
una ventana de longitud $m$ en un algoritmo es una \emph{decisión de modelado}.
Por ejemplo, si $X_1,X_2$ son monedas justas independientes y $X_3=X_1$,
el modelo de orden uno pierde información relevante para el último paso.
\section{Número de parámetros de una tabla condicional}
Para $K$ símbolos y una historia de longitud $m$ (por tanto, $K^m$ contextos posibles), cada contexto posee $K-1$
grado de libertad; la tabla completa usa $K^m(K-1)$ parámetros libres.
Este conteo presupone que cada contexto se trata independientemente, que no
hay simetrías ni parámetros compartidos, y que todas las filas pueden variar.
Las redes neuronales comparten parámetros y no se cuentan de esta forma.

\chapter{Estimación, concavidad y modelos predictivos}
\section{Existencia del estimador de máxima verosimilitud}
Sea $c_i\ge0$, $N=\sum c_i>0$ y
$\ell(p)=\sum_{i:c_i>0}c_i\ln p_i$. Es una función cóncava en el
símplex: el Hessiano en coordenadas positivas observadas es diagonal con
entradas $-c_i/p_i^2$. Para maximizar, la desigualdad de Gibbs o el
multiplicador de Lagrange en el soporte observado dan
$\widehat p_i=c_i/N$. A los índices con $c_i=0$ se les asigna masa nula.
Si se impusiera artificialmente $p_i>0$ para todos, el supremo podría no alcanzarse.
\section{Dirichlet y posterior: derivación completa}
Para $\alpha_i>0$ la densidad de $p\sim\mathrm{Dirichlet}(\alpha)$ es
\[
f(p)=\frac{\Gamma(A)}{\prod_i\Gamma(\alpha_i)}
\prod_{i=1}^Kp_i^{\alpha_i-1},\qquad A=\sum_i\alpha_i.
\]
Con conteos $c$, la verosimilitud es proporcional a $\prod p_i^{c_i}$.
Por Bayes la posterior es Dirichlet$(\alpha_1+c_1,\ldots,\alpha_K+c_K)$.
La identidad de momentos de Dirichlet es
$\mathbb E[p_i\mid c]=(\alpha_i+c_i)/(A+N)$.
Se demuestra integrando, pues la integral de $p_i\prod_j p_j^{\beta_j-1}$
con la constante de normalización es $\beta_i/(\sum_j\beta_j)$.
Así, la predicción posterior de un símbolo nuevo es la media posterior;
no coincide necesariamente con la moda posterior.
\section{Suavizado como estimación de compromiso}
Cuando $\alpha_i=\alpha$,
\[
\frac{c_i+\alpha}{N+K\alpha}
=\frac{N}{N+K\alpha}\frac{c_i}{N}+
\frac{K\alpha}{N+K\alpha}\frac1K.
\]
Esto muestra una retracción hacia la distribución uniforme, no una garantía
de buena generalización. Si $K$ crece, la masa agregada aumenta.
\section{Selección y evaluación separadas}
La elección de $\alpha$, $n$ o las fusiones BPE sobre prueba destruye la
interpretación de prueba como evidencia final independiente. La recomendación
es separar por documento antes de ajustar cualquier representación y realizar
validación de hiperparámetros solo sobre el conjunto de validación.

\chapter{Entropía y divergencia KL}
\section{Entropía como esperanza}
Para una distribución finita $p$, la información de $i$ es
$I_p(i)=-\log_2p_i$ y $H_2(p)=\mathbb E_p[I_p(X)]$.
Con base natural se mide en nats. La relación es
$H_2(p)=H_e(p)/\ln2$. La entropía mide incertidumbre bajo $p$;
la entropía cruzada $H(p,q)$ mide pérdida media al codificar datos de $p$
con probabilidades $q$.
\section{Límites y casos de igualdad}
\textbf{Proposición.} Para $K$ categorías,
$0\le H_2(p)\le\log_2K$. La primera desigualdad se obtiene porque
$-\log_2p_i\ge0$ cuando $0<p_i\le1$; la igualdad ocurre para masas
concentradas en un solo evento. La segunda se deduce de
$D_{\rm KL}(p\|u)=\log_2K-H_2(p)\ge0$, donde $u$ es uniforme.
La igualdad ocurre exactamente para $p=u$.
\section{Desigualdad de Gibbs con masas nulas}
\textbf{Teorema.} $D_{\rm KL}(p\|q)\ge0$, con igualdad si y solo si $p=q$.
\textit{Demostración.} Si existe $i$ con $p_i>0=q_i$, la divergencia es
infinita. En el otro caso, usando $\ln z\le z-1$,
\[
-D_{\rm KL}(p\|q)=\sum_{i:p_i>0}p_i\ln\frac{q_i}{p_i}
\le\sum_{i:p_i>0}(q_i-p_i)\le0.
\]
Para la igualdad, la desigualdad logarítmica exige $q_i=p_i$ en todo el
soporte y la última desigualdad exige que $q$ no tenga masa fuera de él.
Concluimos $p=q$.\hfill$\square$
\section{Entropía cruzada y KL}
Si los soportes son compatibles,
\[
H(p,q)=-\sum_ip_i\log q_i
=-\sum_ip_i\log p_i+\sum_ip_i\log\frac{p_i}{q_i}
=H(p)+D_{\rm KL}(p\|q).
\]
Por tanto $q=p$ minimiza la entropía cruzada sobre \emph{todas} las
distribuciones posibles. Si se optimiza una familia restringida $q_\theta$,
se obtiene el mejor elemento de esa familia según el criterio, no
necesariamente la distribución real.
\section{KL no es una métrica}
Para $p=(3/4,1/4)$ y $q=(1/2,1/2)$,
\[
D_{\rm KL}(p\|q)=1-H_2(p)\approx0.188722\text{ bits},
\quad D_{\rm KL}(q\|p)=\tfrac12\log_2(4/3)\approx0.207519.
\]
La falta de simetría basta para concluir que KL no es una métrica.
Es frecuente usar una divergencia simétrica como Jensen--Shannon:
$\mathrm{JS}(p,q)=\tfrac12D_{\rm KL}(p\|m)+\tfrac12D_{\rm KL}(q\|m)$,
con $m=(p+q)/2$. No se necesita JS en la S01, pero permite situar KL
entre otros objetos de la teoría de la información.
\section{KL entre distribuciones sobre secuencias}
Para $p,q$ sobre $X_{1:T}$, con soportes compatibles,
\[
D_{\rm KL}(p(x_{1:T})\|q(x_{1:T}))=
\sum_{t=1}^T\mathbb E_p\left[
D_{\rm KL}\bigl(p(X_t\mid X_{<t})\|q(X_t\mid X_{<t})\bigr)\right].
\]
Se deriva factorizando el cociente de densidades o masas mediante la regla
de la cadena y conmutando la suma finita y la esperanza. La fórmula es una
versión de la regla de la cadena para KL, y enlaza la modelación de
secuencias con las diferencias de pérdida condicional.

\chapter{Perplejidad, codificación y tokenización}
\section{La perplejidad como transformación de la pérdida}
Para una realización observada y $r_t=q(x_t\mid x_{<t})>0$,
\[
\operatorname{PPL}=\exp\left(-\frac1T\sum_{t=1}^T\ln r_t\right)
=\left(\prod_{t=1}^T\frac1{r_t}\right)^{1/T}.
\]
La igualdad se obtiene aplicando propiedades del logaritmo. Si $r_t=0$
para algún evento observado, la pérdida y la perplejidad son infinitas.
Si todos los eventos reciben probabilidad $1/K$, la PPL vale $K$.
\section{Producto de probabilidades y probabilidad inicial}
Una matriz de transición por sí sola solo especifica las probabilidades
condicionales de pasos consecutivos. Para dar una ley sobre cadenas de
longitud $T$ se debe fijar $p(X_1)$ o condicionar explícitamente al símbolo
inicial. Si la prueba \texttt{abca} tiene transición total $1/35$ bajo
$\alpha=1$ y la inicial es uniforme, su probabilidad conjunta es $1/105$.
Este detalle afecta el número de factores $T$ y el denominador de la PPL.
\section{Bits por carácter y unidades de información}
Para el texto $z$ con $N$ caracteres bajo una convención documentada,
\[
\operatorname{BPC}(z;q)=\frac{-\log_2q(\tau(z))}{N}.
\]
Si $\tau$ es una tokenización reversible y la probabilidad cubre el texto
completo y todas las marcas que lo determinan, esta cantidad representa
una tasa ideal por carácter. De otro modo puede faltar información para
reconstruir el texto. El nombre BPC requiere especificar si se cuentan
puntos de código Unicode, bytes o grafemas y cómo se tratan saltos de línea.
\section{Relación entre PPL y BPC bajo una tokenización fijada}
Si la segmentación de $z$ contiene $T$ tokens y $\operatorname{PPL}_q$
se calculó sobre exactamente esos $T$ eventos, entonces
\[
\operatorname{BPC}=\frac{T}{N}\log_2\operatorname{PPL}_q.
\]
La fórmula es algebraica; no autoriza mezclar valores de PPL, $T$ y $N$
provenientes de convenciones o textos distintos.
\section{Byte Pair Encoding: reducción local}
Sea un corpus representado por secuencias finitas y un par $(u,v)$.
Cada reemplazo \emph{no solapado} de $u,v$ por $uv$ disminuye
el número de símbolos en una unidad. Si se ejecutan $r$ reemplazos,
la longitud total desciende exactamente $r$. Para el par $(a,a)$ en
$\texttt{aaaaa}$ se cuentan cuatro apariciones solapadas pero solo dos
reemplazos no solapados de izquierda a derecha: la disminución es dos.
Por ello, una implementación que cuenta pares brutos debe separar
``frecuencia'' de ``reemplazos efectivos''.
\section{Por qué BPE no prueba un óptimo global}
El algoritmo elige, en cada paso, el par de mayor frecuencia y toma una
decisión local. Aunque disminuye la longitud inmediatamente, cambiar las
fusiones tempranas cambia el conjunto de pares futuros; por ello la
propiedad local no implica optimalidad global con un presupuesto dado de
fusiones. También debe documentarse la política de desempate.
\section{Preguntas suplementarias de nivel matemático}
\begin{enumerate}
\item Pruebe la concavidad de $H(p)$ usando el Hessiano restringido al
interior del símplex.
\item Derive la regla de la cadena de KL a partir de las factorizaciones
$p(x_{1:T})$ y $q(x_{1:T})$.
\item Estudie el límite $\alpha\to\infty$ del estimador suavizado
con $K$ fijo, e interprete el resultado.
\item Dé un ejemplo donde dos modelos tengan igual pérdida en bits por
carácter y perplejidades por token distintas.
\item Explique por qué la validación aleatoria por caracteres es
incorrecta cuando el corpus contiene documentos muy similares.
\end{enumerate}
\clearpage
\chapter*{Dos propiedades que conectan información y modelado}
\addcontentsline{toc}{chapter}{Propiedades que conectan información y modelado}

\section*{La información no aumenta al agrupar categorías}
\addcontentsline{toc}{section}{Agregación y divergencia KL}
Supongamos que un experimento observa tres categorías, pero un instrumento
solo informa si ocurrió $a$ o si ocurrió alguna de las otras dos.
La transformación $f: \{a,b,c\}\to\{a,\mathrm{otro}\}$ agrupa
los resultados $b$ y $c$. Para
$p=(1/2,1/4,1/4)$ y $q=(1/4,1/2,1/4)$, teníamos
$D_{\rm KL}(p\|q)=1/4$ bit. Después de la agregación,
\[
 p_f=(1/2,1/2),\qquad q_f=(1/4,3/4),
\]
y la divergencia entre las distribuciones observables es
\[
\begin{split}
D_{\rm KL}(p_f\|q_f)
&=\frac12\log_2\frac{1/2}{1/4}
+\frac12\log_2\frac{1/2}{3/4}\\
&=\frac12+\frac12\log_2\frac23
\approx0.207519\ \text{bits}.
\end{split}
\]
Ha disminuido respecto del $0.25$ original. Este ejemplo particular
ilustra una propiedad general de la KL denominada \emph{procesamiento de
datos}: al aplicar una transformación determinista común a los
resultados, no aumenta la divergencia.

\begin{figure}[H]\centering
\begin{tikzpicture}[>=Stealth,font=\small,
 node/.style={draw=azul,rounded corners,fill=azulclaro,minimum width=17mm,minimum height=8mm}]
\node[node](a)at(0,1){$a$};\node[node](b)at(0,0){$b$};\node[node](c)at(0,-1){$c$};
\node[node](yes)at(6,0.8){$a$};\node[node](oth)at(6,-0.8){otro};
\draw[->] (a)--(yes);\draw[->] (b)--(oth);\draw[->] (c)--(oth);
\node at (3,-1.8) {La observación final olvida la distinción entre $b$ y $c$.};
\end{tikzpicture}
\caption{La agregación pierde información sobre la diferencia entre dos leyes.}
\end{figure}

\textbf{Justificación.} Si $A_y=f^{-1}(y)$, la desigualdad de suma logarítmica
permite escribir
\[
\sum_{i\in A_y}p_i\ln\frac{p_i}{q_i}
\ge\left(\sum_{i\in A_y}p_i\right)
\ln\frac{\sum_{i\in A_y}p_i}{\sum_{i\in A_y}q_i}.
\]
Sumando en las categorías agregadas se obtiene
$D_{\rm KL}(p\|q)\ge D_{\rm KL}(p_f\|q_f)$.
La prueba exige las convenciones habituales de soporte, y no afirma
que ambas divergencias tengan que ser diferentes: podrían coincidir.
Este resultado es útil en investigación aplicada cuando un reporte
técnico se reduce a etiquetas de eventos; esa reducción puede eliminar
diferencias que existían a nivel textual.

\section*{Por qué ampliar el contexto puede reducir la incertidumbre óptima}
\addcontentsline{toc}{section}{Condicionar reduce la entropía en promedio}
Sea $X$ el próximo símbolo, $Y$ el último símbolo observado y $Z$
una porción anterior del contexto. La teoría de la información
establece
\[
H(X\mid Y,Z)\le H(X\mid Y).
\]
No es una afirmación acerca de que \emph{cualquier modelo entrenado}
con más contexto rendirá mejor. Es una desigualdad sobre las verdaderas
entropías condicionales de una distribución fija.

Para justificarla, definamos la información mutua condicional:
\[
I(X;Z\mid Y)
=\mathbb E_{Y,Z}\Bigl[D_{\rm KL}\bigl(
 p(X\mid Y,Z)\,\|\,p(X\mid Y)
\bigr)\Bigr].
\]
Por Gibbs, el integrando es no negativo en cada contexto del soporte.
Al expandir los logaritmos también se obtiene
\[
I(X;Z\mid Y)=H(X\mid Y)-H(X\mid Y,Z).
\]
Combinando ambas igualdades deducimos la desigualdad.
La igualdad ocurre si $X$ y $Z$ son condicionalmente independientes dado $Y$.

\begin{ejemplo}[Una dependencia que el último símbolo no explica]
En la construcción $X_1$ y $X_2$ Bernoulli$(1/2)$ independientes y
$X_3=X_1$, se tiene
\[
H(X_3\mid X_2)=1\ \text{bit},
\qquad H(X_3\mid X_1,X_2)=0.
\]
La historia más amplia elimina toda incertidumbre \emph{verdadera}
sobre $X_3$, mientras que el bigrama conserva un bit de incertidumbre.
En muestras pequeñas, sin embargo, estimar más contextos puede aumentar
la varianza y empeorar la predicción en prueba. Allí reaparece el
compromiso entre capacidad del modelo, datos y suavizado.
\end{ejemplo}
La transición de este resultado a los Transformers es conceptual:
el acceso a contexto amplio \emph{permite} representar dependencias
que los bigramas descartan, pero aprenderlas de forma fiable exige
datos, parametrización, optimización y evaluación, asuntos de S02--S05.

\chapter*{Profundización: cinco conexiones que preparan la investigación}
\addcontentsline{toc}{chapter}{Cinco conexiones para investigación}

\section*{1. ¿Qué se puede concluir de una cadena de Markov?}
\addcontentsline{toc}{section}{Cadenas de Markov y distribución estacionaria}
El modelo suavizado del ejemplo principal puede estudiarse como una cadena
homogénea de tres estados, con distribución inicial $\pi_0$ y matriz
\[
P=\begin{pmatrix}1/7&3/7&3/7\\1/2&1/3&1/6\\2/5&2/5&1/5\end{pmatrix}.
\]
Por construcción, $P_{ij}\ge0$ y $\sum_jP_{ij}=1$; cada fila define
una probabilidad. Como todas las entradas son estrictamente positivas,
la cadena es irreducible y aperiódica: desde cualquier estado se llega
a cualquier otro en un paso, y cada estado admite una transición a sí mismo.
Bajo estas condiciones existe una única distribución estacionaria $\pi$
con componentes estrictamente positivas y $\pi P=\pi$.

Para calcularla, escribimos $\pi=(u,v,w)$ y resolvemos
\[
\begin{cases}
 u=\tfrac17u+\tfrac12v+\tfrac25w,\\
 v=\tfrac37u+\tfrac13v+\tfrac25w,\\
 w=\tfrac37u+\tfrac16v+\tfrac15w,\\
 u+v+w=1.
\end{cases}
\]
Una de las primeras tres ecuaciones es redundante: las filas de $P$
suman uno. Tras sustituir y resolver el sistema lineal queda
\[
\boxed{\pi=\left(\frac{98}{281},\frac{108}{281},\frac{75}{281}\right).}
\]
Se verifica, por ejemplo, que
\[
\frac{98}{281}\frac17+\frac{108}{281}\frac12
+\frac{75}{281}\frac25
=\frac{14+54+30}{281}=\frac{98}{281}.
\]
El mismo procedimiento verifica las otras dos coordenadas.
Este cálculo revela la distinción entre \emph{probabilidad de transición}
y \emph{probabilidad de ocupación a largo plazo}: no son iguales por
definición. Una distribución inicial arbitraria evoluciona como
$\pi_t=\pi_0P^t$ y se aproxima a la estacionaria bajo las hipótesis
mencionadas. En lenguaje real, las probabilidades condicionadas suelen
cambiar con el contexto y no existe necesariamente una matriz homogénea
fija en ese sentido.

\begin{puente}
Este es el punto de contacto directo con el curso paralelo de Series
de Tiempo: cadenas de Markov, matrices estocásticas, distribución invariante
y pronóstico condicional. Más adelante podremos comparar esta representación
con modelos ocultos de Markov y modelos de estado, pero en S01 no es
necesario introducir sus algoritmos.
\end{puente}

\section*{2. Entropía condicional y tasa de información de una cadena}
\addcontentsline{toc}{section}{Entropía condicional y tasa de información}
La entropía conjunta de dos símbolos satisface
\[
H(X_1,X_2)=H(X_1)+H(X_2\mid X_1).
\]
Esto se demuestra al expandir
$-\sum_{i,j}p(i,j)\log_2 p(i,j)$ y usar
$p(i,j)=p(i)p(j\mid i)$:
\[
\begin{split}
H(X_1,X_2)&=-\sum_{i,j}p(i,j)\log_2p(i)\\
&\quad-\sum_{i,j}p(i,j)\log_2p(j\mid i)\\
&=H(X_1)+\sum_i p(i)H(X_2\mid X_1=i).
\end{split}
\]
El promedio $\sum_i p(i)H(X_2\mid X_1=i)$ es
la entropía condicional: mide cuánto falta por conocer sobre $X_2$
cuando se conoce el estado precedente.

Para la cadena estacionaria construida antes, su tasa de entropía es
\[
h(P,\pi)=-\sum_{i\in V}\pi_i\sum_{j\in V}P_{ij}\log_2P_{ij}.
\]
La fórmula surge porque el símbolo inicial contribuye un término finito
$H(X_1)$, mientras que cada transición, bajo estacionariedad, tiene
la misma entropía condicional media. En efecto,
\[
H(X_{1:T})=H(X_1)+\sum_{t=2}^T H(X_t\mid X_{t-1})
=H(\pi)+(T-1)h(P,\pi).
\]
Dividiendo por $T$ y dejando $T\to\infty$ se obtiene
$\lim_{T\to\infty}T^{-1}H(X_{1:T})=h(P,\pi)$.
Hay dos lecciones: el costo de representar el símbolo inicial
no desaparece en una cadena corta y la entropía marginal $H(\pi)$
no coincide necesariamente con la entropía condicional media.

\begin{ejemplo}[Un contraste con números sencillos]
Si las transiciones son independientes del estado anterior y cada fila
es $(1/2,1/2)$, entonces $H(X_t\mid X_{t-1})=1$ bit y la cadena
no aprovecha memoria. En cambio, para
$P=\left(\begin{smallmatrix}0.9&0.1\\0.1&0.9\end{smallmatrix}\right)$
la estacionaria es uniforme, $H(X_t)=1$ bit, pero
\[
H(X_t\mid X_{t-1})=-0.9\log_20.9-0.1\log_20.1\approx0.4690
\text{ bits}.
\]
El conocimiento del estado anterior reduce la incertidumbre de
la siguiente observación sin cambiar la distribución marginal.
\end{ejemplo}

\section*{3. Soporte, símbolos desconocidos y final de secuencia}
\addcontentsline{toc}{section}{Soporte y símbolos desconocidos}
El suavizado aditivo distribuye masa positiva entre los $K$ símbolos
incluidos en el vocabulario del modelo. \emph{No asigna automáticamente}
probabilidad a cualquier símbolo que no pertenezca a ese vocabulario.
En un modelo de palabras entrenado con $V$, una palabra nueva necesita
una decisión de representación: sustituirse por un símbolo
$\langle\mathrm{UNK}\rangle$, segmentarse en unidades menores o
admitirse mediante algún mecanismo adicional.

\begin{figure}[H]\centering
\begin{tikzpicture}[>=Stealth,font=\small,
 box/.style={draw=azul,fill=azulclaro,rounded corners,align=center,text width=3.1cm,minimum height=9mm}]
\node[box](a)at(0,0){Palabra de prueba\\\texttt{revestidor}};
\node[box](b)at(4.1,1){Vocabulario de palabras\\\texttt{<UNK>}};
\node[box](c)at(4.1,-1){Vocabulario subléxico\\varias unidades};
\draw[->,thick] (a)--(b);\draw[->,thick] (a)--(c);
\end{tikzpicture}
\caption{Dos políticas distintas para un símbolo no visto; no son intercambiables.}
\end{figure}

Si además la longitud de la cadena es variable, necesitamos fijar
cómo termina la secuencia. Una construcción usual introduce un símbolo
de fin $\langle/\mathrm{s}\rangle$ y asigna probabilidades a las
transiciones hasta encontrarlo. De este modo, bajo hipótesis de
terminación adecuadas, el modelo determina una ley sobre cadenas
finitas de longitudes variables. Si se omite la probabilidad de parar,
la suma de probabilidades de cadenas completas puede no estar
bien definida.

\begin{ejemplo}[Parada con probabilidad geométrica]
Considérese un generador que emite un símbolo ordinario con probabilidad
$1-r$ y termina con probabilidad $r\in(0,1)$ en cada paso, con decisiones
independientes. Si $N$ es el número de símbolos ordinarios previos al fin,
\[
P(N=n)=(1-r)^nr,\quad n=0,1,2,\ldots,
\qquad \mathbb E[N]=\frac{1-r}{r}.
\]
Al limitar la generación a un máximo de $m$ símbolos mediante un arnés,
la probabilidad de truncamiento es $P(N\ge m)=(1-r)^m$.
Esta cuenta es una idealización explícita: en un LLM real la probabilidad
de fin depende de la historia y la política de decodificación.
\end{ejemplo}

\section*{4. Códigos ideales y límites de interpretación}
\addcontentsline{toc}{section}{Código ideal y pérdida logarítmica}
La cantidad $-\log_2q(x)$ se interpreta como longitud de código ideal
según $q$. Esto motiva tanto la entropía como la divergencia KL:
si los datos siguen $p$ pero se utiliza $q$,
\[
\mathbb E_p[-\log_2q(X)]-\mathbb E_p[-\log_2p(X)]
=D_{\rm KL}(p\|q).
\]
Es el exceso esperado de longitud ideal, no una promesa de que
cualquier compresor real tenga exactamente esa longitud en un archivo.
Los códigos binarios de longitud entera deben cumplir restricciones
como la desigualdad de Kraft y agregan sobrecostos de cabecera,
fronteras y representación.

\begin{ejemplo}[Cuando la entropía no es el costo de un símbolo concreto]
Para $p=(1/2,1/4,1/4)$, el símbolo de probabilidad $1/2$ contribuye
$1$ bit si ocurre, mientras que cada símbolo de probabilidad $1/4$
contribuye $2$ bits. La entropía $1.5$ bits es un \emph{promedio esperado},
no la longitud de cada realización. Con un código de prefijo
$\{0,10,11\}$ las longitudes individuales son exactamente $(1,2,2)$,
y el promedio coincide con 1.5 bits.
\end{ejemplo}

La pérdida logarítmica es una regla de puntuación propia: su riesgo
esperado se minimiza al informar la distribución verdadera.
La demostración ya está contenida en la identidad de entropía cruzada
más Gibbs. Por ello tiene sentido usarla para evaluar predicciones
probabilísticas, pero no sustituye métricas de exactitud, fidelidad o
utilidad operacional cuando esas son las preguntas científicas.

\section*{5. Incertidumbre experimental y unidad de remuestreo}
\addcontentsline{toc}{section}{Unidad estadística de comparación}
Supongamos que se quiere comparar dos procedimientos A y B en $m$
documentos reservados. Si $L_{A,j}$ y $L_{B,j}$ denotan las pérdidas
en el mismo documento $j$, definimos
\[
d_j=L_{A,j}-L_{B,j},\qquad \overline d=\frac1m\sum_{j=1}^m d_j.
\]
Un análisis pareado retiene la estructura de comparación: cada documento
aporta una diferencia. Si las unidades documentales son razonablemente
independientes, su error estándar clásico es $s_d/\sqrt m$, donde
\[
s_d^2=\frac1{m-1}\sum_{j=1}^m(d_j-\overline d)^2.
\]
Este error estándar y los intervalos que se construyen con él no
resultan válidos por el mero hecho de tener millones de tokens: los
símbolos dentro de un mismo documento comparten dependencias.

En investigación con textos técnicos, varios documentos del mismo pozo,
incidente o empresa también pueden estar correlacionados. En ese caso la
unidad de remuestreo puede necesitar elevarse hasta el pozo, incidente o
fuente, según la pregunta y el proceso de muestreo. Una comparación que
ignora esto puede aparentar una precisión mayor de la real.

\begin{ejemplo}[Diferencias pareadas construidas para ilustrar el método]
Considérense cuatro documentos con diferencias en BPC
$(-0.20,0.10,-0.10,0.00)$, calculadas como A menos B.
Su promedio es $-0.05$ BPC. Los cuadrados respecto de la media son
$0.0225,0.0225,0.0025,0.0025$; la suma es $0.05$.
La varianza muestral es $0.05/3\approx0.01667$ y el error estándar
$\sqrt{0.01667/4}\approx0.06455$. Con solo cuatro documentos,
el valor puntual no debe venderse como efecto establecido. El ejemplo es
\textbf{ILUSTRATIVO}, no procede de una prueba real.
\end{ejemplo}

\section*{6. Cuidado con una supuesta ventaja metodológica}
\addcontentsline{toc}{section}{Separar hipótesis de evidencia}
Un tokenizador puede producir menos unidades por palabra sin que
su modelo probabilístico correspondiente alcance menor pérdida de
predicción; la arquitectura, la cantidad de datos y el tratamiento
de caracteres pueden alterar esa relación. Por eso conviene separar
métricas \emph{descriptivas de segmentación} de métricas \emph{predictivas}.
Si comparamos dos modelos completos, podremos describir diferencias
de tasa; si queremos atribuir la diferencia a BPE, habrá que diseñar un
experimento en el que los otros factores se mantengan comparables.

Esta cautela también explica por qué la matriz IC10 es una \emph{semilla}
de investigación, no un resultado preestablecido. Un resultado nulo,
negativo o contrario a la hipótesis puede ser científicamente útil si
la medición fue defendible. El paso siguiente, en S02, será revisar
candidatos y disponibilidad de datos antes de elegir un problema.

